% ===== PRÉAMBULE CANONIQUE — à coller VERBATIM en tête de CHAQUE .tex =====
\documentclass[11pt,a4paper]{article}
\usepackage[utf8]{inputenc}\usepackage[T1]{fontenc}\usepackage{lmodern}
\usepackage[french]{babel}
\usepackage{amsmath,amssymb,mathtools}\usepackage{icomma}
\usepackage{siunitx}\sisetup{locale=FR}  % \qty{}{}, \si{}, \ang{}
\usepackage{xcolor}\usepackage{colortbl}
\usepackage{tikz}\usepackage{pgfplots}\pgfplotsset{compat=1.18}
\usepackage[siunitx,europeanresistors]{circuitikz} % circuits (élec.)
\usepackage[most]{tcolorbox}\usepackage{enumitem}\usepackage{array,booktabs}
\usepackage[a4paper,margin=2.2cm]{geometry}
\usetikzlibrary{arrows.meta,calc,angles,quotes,positioning,patterns,%
  backgrounds,shapes.geometric,decorations.pathmorphing,decorations.markings,intersections}
\definecolor{primary}{RGB}{222,49,99}\definecolor{primarydark}{RGB}{150,22,55}
\definecolor{defcol}{RGB}{0,114,178}\definecolor{methcol}{RGB}{52,92,148}
\definecolor{excol}{RGB}{0,135,90}\definecolor{attncol}{RGB}{200,40,55}
\tcbset{boxbase/.style={enhanced,breakable,boxrule=0.9pt,arc=2.5pt,
  left=3mm,right=3mm,top=2.3mm,bottom=2.3mm,fonttitle=\bfseries,coltitle=white}}
% --- boîtes TUTORIEL (noms EXACTS, ne pas renommer) ---
\newtcolorbox{cle}[1][]{boxbase,colback=defcol!6,colframe=defcol,
  title={\textsc{À retenir}\ifx\relax#1\relax\relax\else\textmd{ : \itshape #1}\fi}}
\newtcolorbox{methode}[1][]{boxbase,colback=methcol!7,colframe=methcol,sharp corners,
  title={\textsc{Méthode}\ifx\relax#1\relax\relax\else\textmd{ --- \itshape #1}\fi}}
\newtcolorbox{exemplebox}[1][]{boxbase,colback=excol!5,colframe=excol,
  title={\textsc{Exemple}\ifx\relax#1\relax\relax\else\textmd{ : \itshape #1}\fi}}
\newtcolorbox{piege}[1][]{boxbase,colback=attncol!6,colframe=attncol,
  title={\textsc{Piège}\ifx\relax#1\relax\relax\else\textmd{ : \itshape #1}\fi}}
\newtcolorbox{remarque}[1][]{boxbase,colback=black!4,colframe=black!45,
  title={\textsc{Remarque}\ifx\relax#1\relax\relax\else\textmd{ : \itshape #1}\fi}}
% --- boîtes EXERCICE / CORRIGÉ (noms EXACTS, auto-comptées) ---
\newtcolorbox[auto counter]{exo}[1][]{boxbase,colback=primary!4,colframe=primary,
  title={\textsc{Exercice}~\thetcbcounter\ifx\relax#1\relax\relax\else\textmd{ --- \itshape #1}\fi}}
\newtcolorbox[auto counter]{corrige}[1][]{boxbase,colback=excol!4,colframe=excol,
  title={\textsc{Corrigé}~\thetcbcounter\ifx\relax#1\relax\relax\else\textmd{ --- \itshape #1}\fi}}
\newcommand{\vv}[1]{\vec{#1}}\newcommand{\ds}{\displaystyle}
\newcommand{\R}{\mathbb{R}}\newcommand{\N}{\mathbb{N}}\newcommand{\Z}{\mathbb{Z}}
% =========================================================================
\begin{document}
\sisetup{per-mode=symbol}

\begin{corrige}[la pierre au bout d'un fil]
\begin{enumerate}[label=\alph*)]
  \item $v=\dfrac{2\pi R}{T}=\dfrac{2\times\num{3,14}\times\num{0,5}}{1}\approx\qty{3,14}{\metre\per\second}$.
  \item $a_c=\dfrac{v^2}{R}=\dfrac{(\num{3,14})^2}{\num{0,5}}\approx\qty{19,7}{\metre\per\second\squared}$.
  \item $F_c=m\,a_c=\num{0,2}\times\num{19,7}\approx\qty{3,9}{\newton}$.
\end{enumerate}
\end{corrige}

\begin{corrige}[la force centripète]
$F_c=\dfrac{mv^2}{R}=\dfrac{2\times 6^2}{3}=\dfrac{2\times36}{3}=\dfrac{72}{3}=\qty{24}{\newton}$.
\end{corrige}

\begin{corrige}[deux enfants sur un manège]
\begin{enumerate}[label=\alph*)]
  \item $v_1=\omega R_1=2\times1=\qty{2}{\metre\per\second}$ ; $v_2=\omega R_2=2\times3=\qty{6}{\metre\per\second}$.
  \item L'enfant du \emph{bord} ($R_2$) va le plus vite. Ils ont pourtant la \textbf{même} vitesse angulaire
        $\omega$ : c'est $v=\omega R$ qui augmente avec le rayon.
\end{enumerate}
\end{corrige}

\begin{corrige}[retrouver la période]
\begin{enumerate}[label=\alph*)]
  \item $T=\dfrac{2\pi R}{v}=\dfrac{2\times\num{3,14}\times2}{4}=\dfrac{\num{12,56}}{4}\approx\qty{3,14}{\second}$.
  \item $f=\dfrac{1}{T}=\dfrac{1}{\num{3,14}}\approx\qty{0,32}{\hertz}$.
\end{enumerate}
\end{corrige}

\begin{corrige}[doubler la vitesse]
\begin{enumerate}[label=\alph*)]
  \item $F_c=\dfrac{mv^2}{R}\propto v^2$ : si $v$ double, $F_c$ est multipliée par $2^2=\mathbf{4}$.
  \item $F_c'=4\times10=\qty{40}{\newton}$.
\end{enumerate}
\end{corrige}
\end{document}
